%! Author = sriadityavedantam
%! Date = 3/29/26

\section{Division}

Now that we have learned the axioms of arithmetic, let us learn about the division algorithm.

We have all, hopefully, learned how to divide in grade school.
As a revision, you can divide a number evenly by some other number and whatever is left over will result as the remainder.
This can be written more formally as
\[
    \text{dividend} = (\text{divisor})(\text{quotient}) + (\text{remainder}).
\]
Now there is an important understanding I wanted to show to the audience.
Every basic arithmetic operation can be written in terms of addition and multiplication.
We will later see with rings that we make our lives easier by doing subtraction, which shows both an inverse and additive property.
But for now, that is all mumbo-jumbo.
\begin{theorem}[Division Algorithm]{
    Suppose $a,b\in\mathbb{Z}$ with $b>0$.
    Then there exist $q,r\in\mathbb{Z}$ such that
    \[
        a = qb + r
    \]
    with
    \[
        0\leq r < b.
    \]
}
    Let there be a set $S$ such that
    \[
        S \coloneqq \{a-xb : a-xb\geq 0,\ x\in\mathbb{Z}\}.
    \]
    We first check that $S\neq \varnothing$.
    Given $a$ and $b$, find $x$ such that $a-xb\geq 0$.
    If $a\geq 0$, let $x = 0$.
    Then $a-xb = a\geq 0$.
    If $a<0$, let $x = a$.
    Then
    \[
        a - ab = a(1 - b),
    \]
    and since $b>0$, we have $b\geq 1$, therefore $1-b\leq 0$.
    Since $a<0$ as well, it follows that $a(1-b)\geq 0$.
    Since $S\neq \varnothing$, then $S$ is well-ordered.
    Thus there exists $r\in S$ such that $r$ is the smallest element of $S$.
    Claim:
    $r\geq 0$ and $r < b$.
    Since $r\in S$, there exists $q\in\mathbb{Z}$ such that $r = a-qb$ and $r\geq 0$.
    It remains to prove that $r < b$.
    Suppose $r\geq b$.
    Then let
    \begin{align*}
        d &= a - (q + 1)b\\
        &= a - qb - b\\
        &= r - b.
    \end{align*}
    Since $r\geq b$, we have $r-b\geq 0$.
    Thus $d\in S$.
    But $d = r-b < r$, which contradicts the fact that $r$ is the smallest element of $S$.
    Therefore $r < b$.
    Hence
    \[
        a = qb + r
    \]
    with $0\leq r < b$.
\end{theorem}

There is a lot to dissect here.
I want to dedicate special focus to this theorem.
This will lay the foundation, so glance your eyes on this beauty and take it in its glory.
But in all seriousness, this is a really important topic to take in, so let us explain it thoroughly.
Similar to what we have in Figure 2.1 with the dividend equation, we just broke it down and generalized it using proof notation.
So given that $a$ is some dividend, we have divisor $b$, and quotient $q$ that are multiplied, then added with remainder $r$.
There is also a reason why the division algorithm requires that $r$ be less than $b$ but at minimum $0$.
This may be trivial, but if $r$ is greater than $b$, we can subtract $b$ from $r$ and get a new remainder.
It has the most optimized equation.
Now that we understand what we are doing in more understandable terms, let us look at our proof itself and implement it as a core memory as how a child may remember their guardian.

\begin{example}
    Let $S$ be a set of remainders.
    We can do this through example.
    If
    \begin{gather*}
        a = 81,\\
        b = 8,\\
    \end{gather*}
    and $x$ is a variable, then
    \[
        r = a-bx.
    \]
    If we let $x = 1$, then $r = 73$.

    If we let $x = 4$, then $r = 49$.

    If we let $x = 10$, then $r = 1$.

    If we let $x = 11$, then $r = -7$.

    However, $r$ can only be at minimum $0$, therefore $r$ cannot be $-7$.
    Therefore our most optimized $r$ is when $x = 10$.
    Of course, $x$ can go in the opposite direction, since we did not bound $\mathbb{Z}$ only to non-negative integers.
\end{example}

Thus we have shown an example of the division algorithm.
Now that we understand the values that set $S$ can contain, even though we have provided proof, we must still prove this through math and generalize it.
And that is exactly what we spend the rest of the proof doing.

We answer questions in this proof such as, what if $a$ is greater than $0$ or less than $0$.
And what happens if $r$ is greater than $b$, which we show cannot happen if $r$ is the smallest non-negative integer of that form.
For example, we could technically have a solution
\begin{gather*}
    a = 200,\\
    b = 2,\\
    x = 10,\\
    r = 180,\\
\end{gather*}
and this fits the equation $a = bx + r$, but it is not the division algorithm remainder because $r$ is not the smallest allowable one.

\begin{proposition}[Uniqueness in the Division Algorithm]{
    The integers $q,r\in\mathbb{Z}$ in the division algorithm are unique.
}
    Given $a\in\mathbb{Z}$ and $b\in\mathbb{Z}$ with $b>0$, suppose
    \[
        a = q_{1}b + r_1
    \]
    such that $q_1,r_1\in\mathbb{Z}$ and $0\leq r_1 < b$.
    Also suppose that
    \[
        a = q_{2}b + r_2
    \]
    such that $q_2,r_2\in\mathbb{Z}$ and $0\leq r_2 < b$.
    Claim:
    $q_1 = q_2$ and $r_1 = r_2$.
    We have
    \begin{align*}
        a &= q_{1}b + r_1\\
        a &= q_{2}b + r_2.
    \end{align*}
    Subtracting gives
    \begin{align*}
        0 &= (q_1 - q_2)b + r_1 - r_2\\
        r_2 - r_1 &= (q_1 - q_2)b.
    \end{align*}
    Since $0\leq r_1,r_2 < b$, we have
    \[
        -b < r_2 - r_1 < b.
    \]
    Therefore
    \[
        -b < (q_1 - q_2)b < b.
    \]
    Since $b>0$, dividing through by $b$ yields
    \[
        -1 < q_1 - q_2 < 1.
    \]
    Since $q_1-q_2\in\mathbb{Z}$, the only possibility is
    \[
        q_1 - q_2 = 0.
    \]
    Therefore $q_1 = q_2$.
    Then
    \begin{align*}
        0 &= (q_1 - q_2)b + r_1 - r_2\\
        0 &= (0)b + r_1 - r_2\\
        0 &= r_1 - r_2.
    \end{align*}
    Thus $r_1 = r_2$.
\end{proposition}

This proposition demonstrates that $q$ and $r$ are unique, and this is really important to show in math when we are proving an algorithm.
Regardless of what $q$ and $r$ are, if they exist, then they are unique, sounding trivial, but as we see the proof is rather less trivial.

This one is a bit more straightforward, therefore there will not be a conceptualizing analysis on this proof.
This is also just further building the proof techniques we have at our arsenal and allowing one to understand the algorithm through and through.

\begin{definition}[Logical Divide]
    Suppose $a,b\in\mathbb{Z}$.
    Let us define the logical divide of $b$ divides $a$ as $b\mid a$.
    If there exists $q\in\mathbb{Z}$ such that
    \[
        a = bq,
    \]
    then $b\mid a$.
    If $b = 0$ and $a \neq 0$, then $b\nmid a$, because $0q = 0$, and $a\neq 0$.
\end{definition}

There is not a strict name for this definition as far as I know, therefore I created a name for it, Logical Divide.
It is the logical notation for the phrase ``$x$ divides $y$'', and it is trivial to Abstract Algebra.

It is slightly different from, say, previous computationally algebraic courses, where one just computes some division and may even end up with a completed or incomplete, rational or not, answer.
Note that up to now we are only sticking with the integers, and this is a really important fact to keep in mind.
Therefore when we say that $2\mid 4$, then we really mean that $4$ is evenly divisible by $2$, but $3$ does not divide $4$, even if we can write it in terms of a decimal.

Another way we can explain this topic is through the division algorithm.
If it does not look similar, we can write $b$ divides $a$ as
\[
    a = bq + r,
\]
where $r = 0$.
Now does this mean that if $a = 0$, then does $0\mid 0$.
Honestly, it is a debated topic in algebra and number theory.
Some may state yes, others may state no.
But what is important is that many people leave it undefined, for the same reason why your calculator cannot divide $0$ by $0$.
Now if $a = 0$ and $b\mid a$, there is an integer $q$ in $\mathbb{Z}$ such that $a = bq$.
This is a proof we will not get into for the sake of saving time and space, but it is a nice practice exercise.
One proof we will be looking at is the following.

\begin{lemma}{
    Assume $b\mid a$ and $b\neq 0$, so $a=bq$ for some $q\in\mathbb{Z}$.
    Then $-b\mid a$.
}
    We have
    \[
        a = (-b)(-q),
    \]
    so $-b\mid a$, since $-q\in\mathbb{Z}$.
    Similarly, $b\mid -a$.
\end{lemma}

This is just a fun fact to rationalize that these four results are possible:
$b\mid a$, $b\mid -a$, $-b\mid a$, and $-b\mid -a$.
Now on a larger note, we must prove transitivity through logical divides.

\begin{lemma}{
    Suppose $a,b,c\in\mathbb{Z}$.
    If $c\mid b$ and $b\mid a$, then $c\mid a$.
}
    There exist $q_1,q_2\in\mathbb{Z}$ such that
    \[
        a = bq_1
    \]
    and
    \[
        b = cq_2.
    \]
    So
    \[
        a = (cq_2)q_1 = c(q_{2}q_1).
    \]
    Since $q_{2}q_1\in\mathbb{Z}$, we have $c\mid a$.
\end{lemma}

One thing to note is that divisibility is not symmetric, which means that if $b\mid a$, it does not follow that $a\mid b$ unless we are in a special case such as $a = \pm b$.
There is a statement that could be said about linear combinations of $a$ and $b$.
If there is an integer $c$ that divides both $a$ and $b$, then for any integers $x$ and $y$, we have
\[
    c\mid xa+yb.
\]
Therefore, $c$ divides any linear combination of $a$ and $b$.
The proof of this is similar to the previous proof before.
The idea is that if you can write $a$ and $b$ in terms of $c$, then the linear combination can also be written in terms of $c$.
Thus showing divisibility.
Try to implement this on your own.
If it has not been noticeable, there is nothing more to learning a course outside of learning the definitions and theorems.

\begin{definition}[Greatest Common Divisor]
    The GCD of $a$ and $b$, written as $\gcd(a,b) = d$, is the integer $d>0$ such that $d\mid a$ and $d\mid b$, and if $c\mid a$ and $c\mid b$, then $c\mid d$.
\end{definition}

The greatest common divisor is a concept that we have learned in grade school.
If we recall, we can write $\gcd(4,6) = 2$, since $2\mid 4$ and $2\mid 6$.

\begin{theorem}[Linear Combinations of GCD]{
    Let $a,b\in\mathbb{Z}$, not both $0$.
    Let there be a set $S$ such that
    \[
        S \coloneqq \{xa + yb : x,y\in\mathbb{Z},\ xa + yb > 0\}.
    \]
    Then $S\neq \varnothing$, and by the Well-Ordering Principle, $S$ has a smallest element called $d$.
    Then $d = \gcd(a,b)$.
    The key statement is that if $d = \gcd(a,b)$, then there exist $x,y\in\mathbb{Z}$ such that
    \[
        d = xa + yb.
    \]
}
    Let $S\neq \varnothing$.
    Then there exists $d\in S$ such that for all $t\in S$, we have $d \leq t$.
    Since $d\in S$, there exist $x,y\in\mathbb{Z}$ such that
    \[
        d = xa + yb.
    \]
    Now our goal is to prove that $d\mid a$.
    Since $d>0$, by the Division Algorithm there exist $q,r\in\mathbb{Z}$ such that
    \[
        a = qd + r
    \]
    with
    \[
        0\leq r < d.
    \]
    Suppose $r > 0$.
    Then
    \begin{align*}
        r &= a - qd\\
        &= a - q(xa + yb)\\
        &= a - qxa - qyb\\
        &= (1 - qx)a - (qy)b.
    \end{align*}
    So $r$ is a linear combination of $a$ and $b$.
    Since $r > 0$ and $r < d$, then $r\in S$, contradicting the assumption that $d$ is the smallest element of $S$.
    Therefore $r = 0$.
    Hence $a = qd$, so $d\mid a$.
    Similarly, we can show $d\mid b$.
    Now suppose $c\mid a$ and $c\mid b$.
    Then $c\mid xa+yb$, which is a linear combination of $a$ and $b$, and this equals $d$.
    Therefore $c\mid d$.
    Hence $d = \gcd(a,b)$.
\end{theorem}

If the gcd of any two integers ever equals $1$, then we say that $a$ and $b$ are relatively prime.
If they are relatively prime, then by the previous theorem, the linear combination will also equal $1$.

\begin{theorem}{
    Suppose $\gcd(a,b) = 1$ and $c\in\mathbb{Z}$ such that $a\mid bc$.
    Then $a\mid c$.
}
    Since $\gcd(a,b) = 1$, there exist $x,y\in\mathbb{Z}$ such that
    \[
        xa + yb = 1.
    \]
    Therefore,
    \begin{align*}
        xa + yb &= 1\\
        cxa + cyb &= c\\
        (cx)a + y(bc) &= c.
    \end{align*}
    Now $a\mid (cx)a$ and by assumption $a\mid bc$, so $a\mid y(bc)$ as well.
    Therefore $a\mid c$.
\end{theorem}

The last thing we will be looking at in this section is the extended gcd algorithm.
The idea behind this is to use the gcd algorithm and then reverse the process in order to find the factors of the linear combination.
This is more of a computational math.

The gcd algorithm can be written in terms of the Division Algorithm and continuing to find the terms that make up the two factors.
An idea of this is using $\gcd(109,26)$.

\[
    109 = 26(4) + 5.
\]

Because $109$ can be split up by $26$ and have a remainder of $5$, this is no different from having a gcd of $(26,5)$.

\[
    26 = 5(5) + 1.
\]

Now because we are left with a remainder of one, and one can go into any number, then $1$ is our final answer for the gcd of $(109,26)$.
This is a way to do the gcd algorithm through division.

But what if we are to set this the other way around.

\[
    1 = 26 - 5(5).
\]

Similar to what we did before, we are shifting all elements in the equation to create the one above.

\begin{gather*}
    1 = 26 - 5(109 - 26(4)).\\
    1 = (-5)(109) + (21)(26).\\
\end{gather*}

Thus we have found the linear combination factors of the equation.